\section{Conclusion}
\label{sec:rewrite-conclusion}

This paper endogenizes investment in AI research, which in turn affects the pace of economic growth. Even under RSI, computing requires resources that compete with consumption and accumulation of physical capital. Algorithms and model weights, by contrast, are non-rival: using them does not deplete them. Following \citet{romer1990}, I use monopoly rents to make research investment profitable. In practice, firms compete to develop algorithms and train models. However, as
\citet{korinekvipra2025} explain, access to better AI can improve the advantage of a leading developer, so an eventual monopoly does not seem to be far-fetched. The prospect of future monopoly rents can help explain intense competition among firms.  %Allowing entry, competing developers, and knowledge diffusion would help determine how competition changes investment, prices, income distribution, and the pace of economic growth.

With an upper bound on AI efficiency, permanently faster wage growth requires sufficient substitution between AI and labor and sufficiently high attainable efficiency. In this AI-dominated regime, output and AI revenue per worker grow faster than wages, so labor's income share vanishes even as wages rise. AI-dominated growth can permanently increase the return on capital, raising questions about
wealth concentration not examined here \citep{bergbuffiezanna2018,mollrachelrestrepo2022}. When labor remains a production bottleneck, the growth rates of output per worker and wages instead return to the exogenous rate of labor-augmenting technological progress. The upper bound on AI efficiency matters for these growth and distributional outcomes: without it, the unit-elasticity case can sustain both growth rates above this exogenous rate while preserving a positive labor share.

RSI is not sufficient for permanently faster output-per-worker growth, nor does such growth require continuing improvements in AI efficiency. AI efficiency can improve even when labor remains a production bottleneck. Under complementarity, output-per-worker growth cannot permanently exceed the rate of labor-augmenting technological progress, even without an upper bound on AI efficiency. With
sufficient substitution and sufficiently high attainable AI efficiency,
capital accumulation and expanding inference compute sustain permanently
faster growth in both output per worker and wages, even as AI efficiency
converges to its upper bound.

The comparison with competition separates the efficiency needed to
sustain AI-dominated growth from the incentives to reach it. Under free
access to AI models and their improvements, competitive suppliers do not
finance RSI. Nevertheless, sufficiently high initial efficiency sustains
AI-dominated growth without monopoly or further research. From a lower
initial efficiency, monopoly rents can finance improvements that allow
an escape from the labor bottleneck. But markups restrict AI use and
raise the efficiency required relative to marginal-cost pricing.
Competition can therefore sustain AI-dominated growth where monopoly
produces more efficient AI but leaves the economy labor-bottlenecked.
Even when both economies are AI-dominated, monopoly delivers faster
growth only if its efficiency gains outweigh this restriction.
Neither market structure nor efficiency alone determines which economy
grows faster. These comparisons do not rank welfare.

Removing monopoly rents after innovation
can undermine the credibility of the incentives offered beforehand.
Patent buyouts could instead reward the developer while opening access
to existing AI technology \citep{kremer1998}. An extension could study
how to value and finance such buyouts while preserving incentives
for further RSI.

Simulations show that early output-per-worker growth can look similar across economies approaching different long-run regimes. The question remains whether human labor will be essential enough for wages to keep pace with output per worker.



% END REVISED CONCLUSION

%The bottleneck mechanism has antecedents in \citet{aghionjonesjones2019}, and rising wages alongside a declining labor share are not themselves new predictions \citep{jones2026future}. In the equilibria I characterize with private investment in RSI, substitutability alone does not suffice: attainable AI efficiency must alsoexceed a threshold that makes capital accumulation sufficiently productive. Above it, higher attainable efficiency raises long-run growth and the interest rate. The wage-growth premium over the exogenous rate equals the output-per-worker growth premium divided by the substitution elasticity, which measures how much of the growth acceleration reaches wages.

%Complementarity prevents output per worker from growing permanently faster than labor-augmenting technological progress. At unit elasticity, however, continued AI improvement can sustain faster output and wage growth with a positive labor share. With substitution and sufficiently strong research returns, the developer's discounted net profit is unbounded even on a fixed finite horizon. This is a failure of the optimization problem, not a proof that an equilibrium reaches a finite-time singularity. Equilibrium existence without an upper bound remains unresolved for the remaining substitution cases.

%The simulations show that moderate early growth can precede either the AI-dominted regime or the labor-bottleneck regime. long-run regime. Research productivity changes both the initial reallocation of resources and the timing of the divergence, not just the speed along a fixed path. It does not change the limiting outcomes in the finite-upper-bound regimes studied here. These are illustrative scenarios, not joint empirical fits to AI revenues, research expenditure, and prices. Limited information on compute allocation also constrains calibration \citep{korinekmckelvey2026}. I therefore place more weight on the conditional long-run results than on the dates of takeoff.

%The abrupt initial adjustment depends on an unanticipated activation of RSI, reversible capital, and compute supplied without installation delays. Gradual adoption and an explicit compute sector with investment frictions would test how much of this adjustment remains when capacity takes time to expand. The model of \citet{erdiletal2025gate} provides a framework for such frictions. Introducing human research tasks and allowing attainable efficiency to evolve would also test the technological assumptions, drawing on the research networks and bottlenecks in \citet{davidsonetal2026}.

%Permanent monopoly restricts the relationship between AI revenue, prices, and research incentives. Together with the production technology, it can imply a substantial AI revenue share even under complementarity; a profitable AI industry need not signal permanently faster growth. Revenue also differs from net profit after inference and research costs. I retain the AI production weight as an illustrative parameter, not an estimate of the industry's size. The concentration and competition mechanisms discussed by \citet{korinekvipra2025} motivate an extension with entry and competing developers to examine how market structure affects private RSI investment.

%The aggregate labor input excludes differences across occupations and the creation of new tasks in which workers retain a comparative advantage. Incorporating the displacement and reinstatement mechanisms of \citet{acemoglurestrepo2019} would allow labor's role to change through task creation and automation, rather than only through relative input use. The present model cannot determine which workers gain or lose during the transition. Nor should its aggregate substitution threshold be interpreted as a restriction on models with different task structures.

%Finally, a declining labor share does not by itself determine household inequality or workers' economic power. The representative household owns capital and the AI developer, so I do not distinguish workers who rely on wages from households that receive most of the profits. As \citet{jones2026future} emphasizes, asset ownership and redistribution matter for who shares in AI-driven prosperity. In \citet{mollrachelrestrepo2022}, automation increases wealth concentration by raising returns to wealth. Heterogeneous ownership and a fiscal sector would allow me to study how capital income and transfers modify the distribution of gains described here through wages and factor shares.

%\enlargethispage{\baselineskip}
%If there is no bound on AI capabilities, the question becomes whether human labor will remain sufficiently essential for wages to keep pace with the rest of the economy.
